\section{Introduction}
\label{sec:introduction}

A sequence of failed constructions can look exhaustive while still mixing
incompatible evidence.  One construction may supply a selective source,
another a nonempty primitive ledger, and a third a determinant, yet no single
object owns the conjunction.  The danger is especially acute when each repair
changes the phase space, the primitive clock, or the operator category.  A
branch-closing argument must therefore answer two questions before it draws a
negative conclusion: what is the exact request universe, and which typed object
owns every failure?

This distinction mirrors a general lesson about no-go results.  Their force is
relative to stated goals, assumptions, frameworks, and mathematical structures;
changing one of those ingredients is an escape pathway rather than another
failure of the original statement \citep{Dardashti2021NoGo}.  Paper 39 turns
that lesson into a deliberately finite audit.  It does not search for a fifth
affine mechanism.  Instead, it treats the content-addressed outcomes of Papers
35--38 as a typed history and asks whether every request admitted by one
explicit encoding has already received an obstruction or a category-exit
classification.

The timing matters.  The four predecessor outcomes were known when the
14-class, 16-token encoding was assembled.  The encoding was then frozen before
the Paper-39 checker.  Hashes therefore establish byte identity for the audit;
they do not show that the universe was selected prospectively, independently of
the outcomes, or completely with respect to every conceivable idea.  Our claim
is relative exhaustion of this retrospective encoding and nothing stronger.

The proof uses two graph granularities.  The \emph{structural spine} has six
nodes and five edges and records the coarse sequence from the Paper-35 object
firewall through the registry handoff.  The \emph{expanded proof DAG} has 22
nodes and 28 edges.  It preserves the branches, source owners, category exits,
and endpoint obstructions needed for a proof.  A total many-to-one projection
links the two graphs through explicit fibers.  The spine is therefore an
executable summary, not a substitute for the proof DAG.

At each in-domain endpoint we evaluate a retrospective six-coordinate
consolidation of pre-existing Route and source fields,
\begin{equation}
  \Good(c)=\Intrinsic(c)\wedge\Rec(c)\wedge\Selective(c)
  \wedge\OwnedDet(c)\wedge\MarkerOK(c)\wedge\Controls(c).
  \label{eq:good-intro}
\end{equation}
The six coordinates require an intrinsic source, primitive recurrence with
repetitions, arithmetic selectivity, same-object determinant ownership, a
compatible marker, and survival of the frozen controls.  The underlying fields
pre-date Paper 39; their conjunction is part of the retrospective encoding and
is not claimed to have been selected independently of the predecessor outcomes.

Our main result says that every one of the 16 recorded requests is classified.
Eight tokens reach endpoints with a nonempty set of failed coordinates in
\cref{eq:good-intro}; eight are explicit exits with an empty failed-coordinate
set.  At the class level this becomes six obstructed classes, six exit-only
classes, and two mixed classes whose canonical request is obstructed while one
enumerated alternative leaves the category.  The separation is essential:
\Exit{} is not evidence for $\neg\Good$.

The finite-graph aspect has close predecessors.  In particular, the Equational
Theories Project fixes a homogeneous finite universe of equational laws and
uses formal proofs and countermodels to complete its implication graph
\citep{BolanEtAl2025ETP}.  Typed proof-dependency graphs and status-bearing
blueprints are also established.  Paper 39 claims neither the first finite
mathematical graph nor the first typed DAG.  Its narrower contribution is the
content-addressed integration of heterogeneous historical repairs, explicit
object/marker/operator exits, and a governance handoff for this affine program.

The contributions are fourfold.
\begin{enumerate}
  \item We freeze a finite theorem domain with exactly 14 repair classes, 16
  request tokens, and 17 internal tags, and we expose both graph granularities
  with an exact $22/28\to6/5$ projection.
  \item We prove endpoint-obstruction totality over the encoded domain: every
  obstructed token has a specified endpoint with a nonempty failed-\Good{} set,
  while every exit token has a typed boundary and no obstruction credit.
  \item We isolate two ownership hazards.  Edge \code{E22} is a non-domain
  historical firewall with empty coverage fibers, and the transition
  \code{E07}/\code{E36\_37} resets object, marker, operator owner, and
  determinant owner under the Paper-37 source lock.
  \item We derive the strict protocol consequence without inflating it into
  candidate success: all Route-A coordinates fail, Route B stays locked, and
  control returns to the historical registry without ranking or selecting an
  entry.
\end{enumerate}

\Cref{fig:spine-expanded} summarizes the exact relationship between the two
graph granularities.

\begin{figure}[t]
  \centering
  \resizebox{\textwidth}{!}{\begin{tikzpicture}[
  x=1cm,y=1cm,>=Latex,
  spine/.style={rounded corners=2pt,draw=formalblue,very thick,fill=white,
    align=center,text width=2.25cm,minimum height=1.15cm,font=\scriptsize},
  close/.style={spine,draw=stopred,text=stopred},
  registry/.style={spine,draw=governancepurple,text=governancepurple},
  fiber/.style={rounded corners=2pt,draw=charcoal,thick,fill=softgray,
    align=left,text width=2.55cm,minimum height=1.62cm,font=\scriptsize,
    inner sep=4pt},
  aux/.style={fiber,draw=warningamber,dashed,fill=white},
  arr/.style={->,very thick,formalblue},
  terminalarr/.style={->,very thick,stopred},
  governarr/.style={->,very thick,governancepurple},
  projection/.style={->,densely dotted,thick,charcoal},
  badge/.style={rounded corners,draw=charcoal,fill=white,align=center,
    font=\scriptsize,inner sep=4pt}
]
\node[spine] (s35) at (0,2.6)
  {\texttt{N35}\\object firewall};
\node[spine] (s36) at (3.0,2.6)
  {\texttt{N36}\\filling / clock};
\node[spine] (s37) at (6.0,2.6)
  {\texttt{N37}\\coefficient fork};
\node[spine] (s38) at (9.0,2.6)
  {\texttt{N38}\\tree / orbital};
\node[close] (s39) at (12.0,2.6)
  {\texttt{N39}\\encoded branch closed};
\node[registry] (sr) at (15.0,2.6)
  {registry handoff\\no ranking};

\draw[arr] (s35) -- (s36);
\draw[arr] (s36) -- (s37);
\draw[arr] (s37) -- (s38);
\draw[terminalarr] (s38) -- (s39);
\draw[governarr] (s39) -- (sr);

\node[badge,above=5mm of s36] (contract)
  {retrospective contract\\\textbf{14 classes / 16 tokens / 17 tags}};
\node[badge,above=5mm of s38,text width=5.1cm] (partition)
  {$28=17$ internal $+5$ closure $+3$ token exits\\
   $+1$ non-domain firewall $+2$ guards};

\node[fiber] (f35) at (0,-0.15)
  {\textbf{P35 fiber (7)}\\
   \texttt{N35F,N35P,N35S,}\\
   \texttt{N35H,N35Q,N35D,}\\
   \texttt{N35B}};
\node[fiber] (f36) at (3.4,-0.15)
  {\textbf{P36 fiber (2)}\\
   \texttt{N36F,N36G}\\[2pt]
   filling, quotient clock};
\node[fiber] (f37) at (6.8,-0.15)
  {\textbf{P37 fiber (3)}\\
   \texttt{N37O,N37D,N37N}\\[2pt]
   local systems, grading};
\node[fiber] (f38) at (10.2,-0.15)
  {\textbf{P38 fiber (5)}\\
   \texttt{N38T,N38M,N38L,}\\
   \texttt{N38O,N38K}};
\node[aux,text width=5.2cm,minimum height=1.9cm,font=\scriptsize] (fa) at (14.0,-0.15)
  {\textbf{exact singleton mappings (5)}\\
   \texttt{N00} $\mapsto$ \texttt{AUX\_CONTRACT\_ROOT}\\
   \texttt{NX} $\mapsto$ \texttt{AUX\_NONMEMBERSHIP\_SINK}\\
   \texttt{NC} $\mapsto$ \texttt{N39\_AFFINE\_BRANCH\_CLOSED}\\
   \texttt{NR} $\mapsto$ \texttt{N\_REGISTRY\_HANDOFF}\\
   \texttt{NS} $\mapsto$ \texttt{AUX\_EMPTY\_REGISTRY\_FALLBACK}};

\draw[projection] (f35.north) -- node[right,font=\tiny]{fiber} (s35.south);
\draw[projection] (f36.north) -- node[right,font=\tiny]{fiber} (s36.south);
\draw[projection] (f37.north) -- node[right,font=\tiny]{fiber} (s37.south);
\draw[projection] (f38.north) -- node[right,font=\tiny]{fiber} (s38.south);

\node[badge,below=4mm of f37,text width=12.2cm] (semantics)
  {total many-to-one projection with explicit, auditable fibers\\
   the 6/5 spine communicates; the 22/28 DAG owns the proof\\
   \texttt{E35\_36} $\cdot$ \texttt{E36\_37} $\cdot$ \texttt{E37\_38}
   $\cdot$ \texttt{E38\_CLOSE} $\cdot$ \texttt{E\_CLOSE\_REGISTRY}};
\end{tikzpicture}
}
  \caption{Two graph granularities encode one relative theorem.  The upper
  six-node/five-edge structural spine summarizes the predecessor sequence and
  governance handoff.  The lower row gives four nontrivial fibers and five
  exact singleton mappings, together retaining all 22 expanded nodes; the 28
  expanded edges are partitioned as shown.  The projection is total and
  many-to-one with explicit fibers, not an isomorphism.  Closure applies only
  to the retrospective 14-class, 16-token, 17-tag encoding.}
  \label{fig:spine-expanded}
\end{figure}

\Cref{sec:related} positions the audit against no-go methodology, proof
certificates, and finite graph completion.  \Cref{sec:contract} defines the
contract and typed graphs.  \Cref{sec:theorem} proves relative exhaustion.
\Cref{sec:ownership} isolates ownership transfer, and
\cref{sec:executable,sec:route} state the executable trust boundary and Route
consequence.  Full token and proof ledgers appear in the appendices.
